% SPDX-FileCopyrightText: 2026 Will Cook
% SPDX-License-Identifier: CC-BY-4.0
\documentclass[11pt]{article}
\usepackage[T1]{fontenc}
\usepackage[utf8]{inputenc}
\usepackage[a4paper,margin=1in]{geometry}
\usepackage{amsmath,amssymb,amsthm,booktabs,seqsplit}
\usepackage[dvipsnames]{xcolor}
\usepackage[colorlinks=true,linkcolor=black,urlcolor=RoyalBlue,citecolor=black]{hyperref}
\newcommand{\commit}{v0.6.0}
\newcommand{\repobase}{https://github.com/wcook04/plectis-lean-erdos249-257/blob/\commit}
\newcommand{\PK}{Erdos249257}
% SPDX-FileCopyrightText: 2026 Will Cook
% SPDX-License-Identifier: CC-BY-4.0
% Generated by scripts/build_paper_module_aliases.py; do not edit.
% Each CamelCase component contributes at most its first three characters.
\newcommand{\modulesigil}[1]{\csname modulesigil@#1\endcsname}
\expandafter\def\csname modulesigil@ActualForeignResidueProjection.lean\endcsname{ActForResPro}
\expandafter\def\csname modulesigil@AdelicHeightObstruction.lean\endcsname{AdeHeiObs}
\expandafter\def\csname modulesigil@BooleanMobiusCarry.lean\endcsname{BooMobCar}
\expandafter\def\csname modulesigil@CampbellShiftSynchronization.lean\endcsname{CamShiSyn}
\expandafter\def\csname modulesigil@CarrySurvivorExtinction.lean\endcsname{CarSurExt}
\expandafter\def\csname modulesigil@CertificateKernel.lean\endcsname{CerKer}
\expandafter\def\csname modulesigil@DiagonalFlexibleOddWindowAffine.lean\endcsname{DiaFleOddWinAff}
\expandafter\def\csname modulesigil@DiagonalFlexibleOddWindowSupply.lean\endcsname{DiaFleOddWinSup}
\expandafter\def\csname modulesigil@DiagonalFreshLossBridge.lean\endcsname{DiaFreLosBri}
\expandafter\def\csname modulesigil@DiagonalPincerCertificateT64Endpoint.lean\endcsname{DiaPinCerT64End}
\expandafter\def\csname modulesigil@DiagonalPincerCertificatesT64.lean\endcsname{DiaPinCerT64}
\expandafter\def\csname modulesigil@DiagonalPincerDecomposition.lean\endcsname{DiaPinDec}
\expandafter\def\csname modulesigil@DiagonalPincerPrimeCertificates/ClosureT64.lean\endcsname{DiaPinPriCerT64}
\expandafter\def\csname modulesigil@FreshPrimeDeficitDecomposition.lean\endcsname{FrePriDefDec}
\expandafter\def\csname modulesigil@FullTargetPrimeAdjunctionNoGo.lean\endcsname{FulTarPriAdjNoGo}
\expandafter\def\csname modulesigil@GapFareyBound.lean\endcsname{GapFarBou}
\expandafter\def\csname modulesigil@GcdMomentCalculus.lean\endcsname{GcdMomCal}
\expandafter\def\csname modulesigil@GenericTailOrbitRigidity.lean\endcsname{GenTaiOrbRig}
\expandafter\def\csname modulesigil@GreedyAchievementSet.lean\endcsname{GreAchSet}
\expandafter\def\csname modulesigil@LcmConeFlatness.lean\endcsname{LcmConFla}
\expandafter\def\csname modulesigil@LcmConeNonflat.lean\endcsname{LcmConNon}
\expandafter\def\csname modulesigil@LcmDiagonalReduction.lean\endcsname{LcmDiaRed}
\expandafter\def\csname modulesigil@MersenneLambertLadder.lean\endcsname{MerLamLad}
\expandafter\def\csname modulesigil@MersenneShadowDenominatorGrowth.lean\endcsname{MerShaDenGro}
\expandafter\def\csname modulesigil@MersenneTailAtoms.lean\endcsname{MerTaiAto}
\expandafter\def\csname modulesigil@MobiusSignSupportNoGo.lean\endcsname{MobSigSupNoGo}
\expandafter\def\csname modulesigil@PowerTwoBitLift.lean\endcsname{PowTwoBitLif}
\expandafter\def\csname modulesigil@PowerTwoCenteredBitLift.lean\endcsname{PowTwoCenBitLif}
\expandafter\def\csname modulesigil@PrimitiveDeterminantLift.lean\endcsname{PriDetLif}
\expandafter\def\csname modulesigil@PrimitiveWeightCertificate.lean\endcsname{PriWeiCer}
\expandafter\def\csname modulesigil@RationalSupportCarrySkeleton.lean\endcsname{RatSupCarSke}
\expandafter\def\csname modulesigil@ResidualGaugeObstruction.lean\endcsname{ResGauObs}
\expandafter\def\csname modulesigil@SignedQMomentObstruction.lean\endcsname{SigQMomObs}
\expandafter\def\csname modulesigil@SquareCRTCube.lean\endcsname{SquCRTCub}
\expandafter\def\csname modulesigil@SquaredMersenneDiagonalEnclosure.lean\endcsname{SquMerDiaEnc}
\expandafter\def\csname modulesigil@SternBrocotRunGeometry.lean\endcsname{SteBroRunGeo}
\expandafter\def\csname modulesigil@SublogDivisorCoverage.lean\endcsname{SubDivCov}
\expandafter\def\csname modulesigil@TotientCarryKernelRigidity.lean\endcsname{TotCarKerRig}
\expandafter\def\csname modulesigil@TotientMahlerDefect.lean\endcsname{TotMahDef}
\expandafter\def\csname modulesigil@TotientTailPeriodKiller.lean\endcsname{TotTaiPerKil}

\newcommand{\idn}[1]{\texttt{\expandafter\seqsplit\expandafter{\detokenize{#1}}}}
\newcommand{\lref}[3]{\href{\repobase/\PK/#1\#L#2}{\texttt{\modulesigil{#1}}}}
\emergencystretch=2em
\newtheorem{theorem}{Theorem}[section]
\newtheorem{proposition}[theorem]{Proposition}
\theoremstyle{remark}\newtheorem*{remark}{Remark}
\newcommand{\Z}{\mathbb Z}
\newcommand{\ph}{\varphi}
\title{Transport and Curvature Certificates for the Totient Constant\\
{\large a Lean~4 companion to the Erd\H{o}s~\#249 formalisation}}
\author{Will Cook}
\date{July 2026}

\begin{document}
\sloppy
\maketitle
\begin{abstract}
Let $S=\sum_{n\geq1}\ph(n)/2^n$ and
$R_N=\sum_{j\geq1}\ph(N+j)/2^j$.  The irrationality of $S$ is
Erd\H{o}s Problem~249 and remains open.  We formalise exact secondary
reductions built from finite differences of the tails $R_N$: curvature,
exponent-transport and prime-jump commutators, and an anchored four-vertex
operator at $H,3H,5H,15H$.  Each object splits into an integer dyadic window
and a rigorously bounded tail, giving a decidable modular certificate for
non-integrality.  We prove that old M\"obius divisor channels are affine and
are annihilated by the relevant moments, isolate the new-channel migration at
a prime jump, and show that every fixed-precision valuation--unit word admits
a centred completion.  The finite consumers and no-go theorem are
unconditional; every irrationality conclusion remains conditional on an
explicitly named unbounded supply.  No such supply, and hence no solution of
Problem~249, is claimed.
\end{abstract}

\section{Scope and logical boundary}\label{sec:scope}
The main exposition~\cite{mainpaper} develops the direct certificate reduction
for $S$, proves an unconditional denominator exclusion, and records the open
boundary.  This companion studies finite-difference operators on the same
tail sequence.  They annihilate affine contributions from divisor channels
already present at an lcm height, but the remaining new channels must still be
shown to produce non-integral values at unboundedly many heights.  Thus every
endpoint below has the form
\[
 \text{unbounded supply of finite certificates}\Longrightarrow S\notin\mathbb Q,
\]
with the antecedent open.  A certificate at one height is finite evidence,
not evidence for the cofinal quantifier.  The fixed-precision no-go theorem is
a boundary on one local strategy, not a global impossibility result.

The Lambert-series background is classical.  The open problem is recorded in
the Erd\H{o}s--Graham collection and maintained catalogue
\cite{erdosgraham1980,erdosproblems}.  We make no priority claim for finite
differences or affine moment cancellation as general techniques.  The
contribution here is their precise integration into checked identities,
finite consumers, conditional sockets, and an obstruction theorem.

Blue source sigils abbreviate each CamelCase module component to at most its
first three letters (for example, \texttt{PrimeJumpWindow} becomes
\texttt{PriJumWin}); every sigil links to the exact declaration in that module.

\section{Dyadic windows}\label{sec:windows}
For a finite integer combination $X(H)=\sum_i c_iR_{m_iH}$, truncate each
tail at depth $L$.  Multiplication by $2^L$ gives
\[
 2^LX(H)=A(H,L)+E(H,L),
\]
where $A(H,L)\in\Z$ is computable.  If $|E(H,L)|\le B(H,L)$, then
$X(H)\notin\Z$ whenever
\[
 B(H,L)<A(H,L)\bmod2^L<2^L-B(H,L).
\]
The reusable soundness theorem is
\idn{notMem_int_of_window_remainder_bound}
\lref{PrimeJumpWindow.lean}{53}{notMem_int_of_window_remainder_bound}.
Lean checks the exact decomposition, radius, and concrete modular arithmetic;
an external analytic or arithmetic argument must provide unbounded instances.

\section{Curvature on the lcm cone}\label{sec:curvature}
Set
\[
 Q_H=R_{4H}-3R_{2H}+2R_H.
\]
Its depth-$L$ window $C(H,L)$ satisfies
\[
 2^LQ_H=C(H,L)+E_C(H,L),\qquad |E_C(H,L)|\le6H+3L+3.
\]
The exact split and bound are
\idn{two_pow_mul_curvature_eq_window_add_shifted}
\lref{CurvatureCarry.lean}{119}{two_pow_mul_curvature_eq_window_add_shifted}
and \idn{abs_curvatureShiftedTail_le}
\lref{CurvatureCarry.lean}{137}{abs_curvatureShiftedTail_le}.
The central residue criterion therefore proves $Q_H\notin\Z$ by
\idn{curvature_notMem_int_of_sharpCurvatureCert}
\lref{CurvatureCarry.lean}{159}{curvature_notMem_int_of_sharpCurvatureCert}.
An unbounded supply implies irrationality through
\idn{irrational_totientSeries_of_sharpCurvatureSupply}
\lref{CurvatureCarry.lean}{206}{irrational_totientSeries_of_sharpCurvatureSupply},
but that supply is not proved.

The window evolves by
$C(H,L+1)=2C(H,L)+d(H,L+1)$, checked in
\idn{curvatureWindow_succ}
\lref{CurvatureCarry.lean}{259}{curvatureWindow_succ}.  Exact phase-certificate
and common-carrier formulas expose possible routes to a supply theorem without
asserting one.

\section{Exponent transport and affine cancellation}\label{sec:transport}
If $p\mid H$, multiplication by $p$ adds no squarefree divisor support:
\[
 d\mid pH\Longleftrightarrow d\mid H\qquad(d\text{ squarefree}).
\]
This is \idn{squarefree_dvd_mul_iff_of_dvd}
\lref{ExponentOnlyTransport.lean}{36}{squarefree_dvd_mul_iff_of_dvd}.
The exact residue kernel is affine, not merely periodic:
\[
 K_d(N)=N\,a_d(N)+b_d(N),
\]
where both $a_d$ and $b_d$ are periodic modulo $d$.  The decomposition is
\idn{transportResidueKernel_eq_affineState}
\lref{ExponentOnlyTransport.lean}{141}{transportResidueKernel_eq_affineState}.
Keeping both blocks prevents an intercept-only matrix calculation from being
misread as global non-cancellation.  The exact $H\mapsto3H$ two-block transfer
is \idn{residueTransportDefect_three_eq_affineState}
\lref{ExponentOnlyTransport.lean}{176}{residueTransportDefect_three_eq_affineState}.

Every old channel $d\mid H$ is killed by the first-order defect
\idn{residueTransportDefect_eq_zero_of_dvd}
\lref{ExponentOnlyTransport.lean}{112}{residueTransportDefect_eq_zero_of_dvd}.
The endpoint
\idn{irrational_totient_series_of_exponentOnlyThreeTransportSupply}
\lref{ExponentOnlyTransport.lean}{205}{irrational_totient_series_of_exponentOnlyThreeTransportSupply}
is conditional on a named fixed-three tower supply.

More generally, coefficients satisfying $\sum_i c_i=\sum_i c_im_i=0$
annihilate every old affine channel.  This is
\idn{oldChannel_affine_moment_annihilation}
\lref{JointExponentTransport.lean}{36}{oldChannel_affine_moment_annihilation}.
For $q(X,Y)=XY-3X-2Y+4$, the vanishing $q(1,1)=q(3,5)=0$ gives
\[
 K_d(15H)-3K_d(3H)-2K_d(5H)+4K_d(H)=0\quad(d\mid H),
\]
checked by \idn{joint35_oldChannel_zero}
\lref{JointExponentTransport.lean}{71}{joint35_oldChannel_zero}.
The associated window radius is $19H+5L+5$; its finite consumer is
\idn{joint35Cone_notMem_int_of_cert}
\lref{JointExponentTransport.lean}{172}{joint35Cone_notMem_int_of_cert}.
The conditional endpoint is
\idn{irrational_totient_series_of_anchoredJoint35ConeSupply}
\lref{JointExponentTransport.lean}{217}{irrational_totient_series_of_anchoredJoint35ConeSupply}.

\section{Prime jumps and three-transport boundaries}\label{sec:primejump}
When $p\nmid H$, a jump from $H$ to $pH$ creates new squarefree support.  The
residue Hecke defect splits into a post-jump channel and an explicit migration
correction in \idn{residueHeckeDefect_eq_postJump_add_migration}
\lref{PrimeJumpMigration.lean}{38}{residueHeckeDefect_eq_postJump_add_migration}.
For $d\mid H$ the total old-channel defect is zero
\lref{PrimeJumpMigration.lean}{46}{residueHeckeDefect_eq_zero_of_dvd}; the
$(d,H,p)=(5,12,5)$ calculation records the two cancelling terms.

Set
\[
 J(H,p)=(R_{2pH}-R_{pH})-p(R_{2H}-R_H).
\]
Its exact remainder radius is $3pH+(p+1)(L+2)$, proved with the window split by
\idn{primeJumpTailCommutator_eq_window_add_remainder}
\lref{PrimeJumpWindow.lean}{91}{primeJumpTailCommutator_eq_window_add_remainder}.
The finite consumer is
\idn{primeJumpTailCommutator_notMem_int_of_sharpKill}
\lref{PrimeJumpWindow.lean}{178}{primeJumpTailCommutator_notMem_int_of_sharpKill}.
Lean verifies $(H,p,L)=(12,5,15)$ in
\idn{primeJumpSharpKill_twelve_five}
\lref{PrimeJumpWindow.lean}{186}{primeJumpSharpKill_twelve_five}; this is one
finite deposit, not an unbounded sequence.

The balanced vertices $H,2H,3H,6H$ give a related radius $9H+4L+4$.
The consumer is \idn{threeTransport_notMem_int_of_cert}
\lref{ThreeTransportBoundary.lean}{126}{threeTransport_notMem_int_of_cert};
the irrationality theorem
\idn{irrational_totientSeries_of_roughnessBoundaryLatticeSupply}
\lref{ThreeTransportBoundary.lean}{177}{irrational_totientSeries_of_roughnessBoundaryLatticeSupply}
is conditional.  The concrete check $(H,L)=(60,12)$ is
\idn{sharpThreeTransportCert_60_12}
\lref{ThreeTransportBoundary.lean}{203}{sharpThreeTransportCert_60_12}.

\section{A first-harmonic interface}\label{sec:harmonic}
For the window discrepancy $A(h,N,L)$ define
\[
 W_1(h,N,L)=\cos\left(2\pi\frac{A(h,N,L)\bmod2^L}{2^L}\right).
\]
Assume $N\in[X,2X)$ and $16(2X+h+L+2)\le2^L$.  If every $N$ failed the
central-residue certificate, each phase would lie within $\pi/8$ of an
integer turn, so each cosine would exceed $9/10$.  Hence
\[
 \sum_{X\le N<2X}W_1(h,N,L)\le\frac9{10}X
\]
forces a certificate in the block.  This implication is
\idn{exists_certifiedKill_of_first_harmonic_gap}
\lref{FirstHarmonicGap.lean}{104}{exists_certifiedKill_of_first_harmonic_gap}.
No cofinal exponential-sum estimate establishing its hypothesis is proved.

\section{Fixed local precision cannot finish the argument}\label{sec:nogo}
A finite valuation--unit symbol records a bounded amount of local transport
data.  For every positive precision $u$, every finite compatible word admits
a prefix-locked centred completion.  Thus no finite word over this fixed local
alphabet excludes all centred endpoints.  The exact statement is
\idn{fixedPrecisionTropicalNoGo}
\lref{TropicalCurvatureCarry.lean}{137}{fixedPrecisionTropicalNoGo}.
This does not say curvature certificates are rare.  A proof may use precision
growing with scale, global correlations, or a harmonic gap.  It rules out only
the inference that a bounded local signature alone forces escape.

\section{Statement matrix and open boundary}\label{sec:index}
\begin{center}\small
\begin{tabular}{@{}p{0.24\textwidth}p{0.46\textwidth}p{0.2\textwidth}@{}}
\toprule Layer & Checked statement & Status \\
\midrule
Curvature & exact split, radius, recurrence, consumer & unconditional \\
Exponent transport & stable support, affine state, cancellation & unconditional \\
Joint $(3,5)$ transport & moment annihilation and consumer & unconditional \\
Prime jump & migration, sharp window, $(12,5,15)$ & finite \\
Three transport & sharp window and $(60,12)$ & finite \\
First harmonic & cosine gap forces a certificate & conditional input \\
Tropical signature & fixed precision has a completion & no-go \\
Irrationality of $S$ & follows from a named unbounded supply & supply open \\
\bottomrule
\end{tabular}
\end{center}
The development proves no cofinal supply of curvature, prime-jump, joint, or
three-transport certificates, and no cofinal first-harmonic gap.  It therefore
does not prove Erd\H{o}s~\#249.

\begin{thebibliography}{9}
\bibitem{erdosgraham1980} P.~Erd\H{o}s and R.~L. Graham,
\emph{Old and New Problems and Results in Combinatorial Number Theory},
Monographies de l'Enseignement Math\'ematique 28 (1980), 61--62,
\href{https://mathweb.ucsd.edu/~ronspubs/80_11_number_theory.pdf}{scan}.
\bibitem{erdosproblems} T.~Bloom (ed.), \emph{Erd\H{o}s Problems}, entry 249,
\href{https://www.erdosproblems.com/249}{\texttt{erdosproblems.com/249}},
accessed July 2026.
\bibitem{mainpaper} W.~Cook, \emph{Erd\H{o}s Problems 249 and 257 in Lean~4:
formalised results, certificate reductions, and open questions}, companion
exposition in this release.
\end{thebibliography}
\end{document}
